\documentclass[a4paper,11pt]{article}
\usepackage[T1]{fontenc}
\usepackage[utf8]{inputenc}
\usepackage[left=2cm,right=2cm,top=2.5cm,bottom=2.5cm]{geometry}
\setlength{\headheight}{14pt}
\usepackage{xcolor}
\usepackage{mdframed}
\usepackage{parskip}
\usepackage{fancyhdr}
\usepackage{hyperref}

\definecolor{usercolor}{RGB}{30, 100, 200}
\definecolor{agentcolor}{RGB}{30, 150, 80}
\definecolor{doccolor}{RGB}{200, 90, 10}
\definecolor{userbg}{RGB}{235, 244, 255}
\definecolor{agentbg}{RGB}{235, 250, 238}
\definecolor{docbg}{RGB}{255, 243, 220}
\definecolor{answerbg}{RGB}{250, 250, 235}

\newmdenv[backgroundcolor=userbg,linecolor=usercolor,linewidth=1.5pt,
  leftmargin=0pt,rightmargin=80pt,innerleftmargin=10pt,innerrightmargin=10pt,
  innertopmargin=8pt,innerbottommargin=8pt,roundcorner=5pt,
  skipabove=6pt,skipbelow=3pt]{userbubble}

\newmdenv[backgroundcolor=agentbg,linecolor=agentcolor,linewidth=1.5pt,
  leftmargin=80pt,rightmargin=0pt,innerleftmargin=10pt,innerrightmargin=10pt,
  innertopmargin=8pt,innerbottommargin=8pt,roundcorner=5pt,
  skipabove=3pt,skipbelow=3pt]{agentbubble}

\newmdenv[backgroundcolor=docbg,linecolor=doccolor,linewidth=1pt,
  leftmargin=80pt,rightmargin=0pt,innerleftmargin=10pt,innerrightmargin=10pt,
  innertopmargin=6pt,innerbottommargin=6pt,roundcorner=4pt,
  skipabove=2pt,skipbelow=3pt]{docenv}

\newmdenv[backgroundcolor=answerbg,linecolor=gray,linewidth=0.5pt,
  leftmargin=80pt,rightmargin=0pt,innerleftmargin=10pt,innerrightmargin=10pt,
  innertopmargin=5pt,innerbottommargin=5pt,roundcorner=4pt,
  skipabove=2pt,skipbelow=6pt]{refanswerenv}

\pagestyle{fancy}
\fancyhf{}
\renewcommand{\headrulewidth}{0.3pt}
\cfoot{\thepage}

\lhead{\textcolor{gray}{\small RAG Conversation Dataset}}
\rhead{\textcolor{gray}{\small Conv.\,\#2\ --\ govt0611}}
\begin{document}

\begin{center}
{\LARGE\bfseries RAG Conversation #9}
\par\medskip
{\small\textcolor{gray}{Domain:\ \texttt{mt-rag-fiqa-beir-elser-512-100-20240501}\quad|\quad Status:\ \textbf{accepted}}}
\end{center}
\vspace{0.4em}
\noindent\textcolor{gray}{\rule{\linewidth}{0.6pt}}
\vspace{0.4em}
\vspace{0.5em}
\noindent{\small\textcolor{usercolor}{\textbf{$\triangleright$\enspace User --- Turn~1}}}\hfill{\tiny\textcolor{gray}{Factoid, Summarization $\cdot$ ANSWERABLE}}
\begin{userbubble}
What approaches are there for pricing a small business?
\end{userbubble}
\noindent\hfill{\small\textcolor{agentcolor}{\textbf{Agent\enspace$\triangleleft$}}}
\begin{agentbubble}
      One approach is to determine the lowest amount a seller would accept, which is the liquidation value. This involves calculating the value of the business's assets, such as land, house, furnishings, and accounts payable, minus any liabilities. Another approach is to price the business based on a multiple of its annual earnings. This can be done using a discounted cash flow analysis, where you project the business's earnings out in perpetuity and discount it by the going interest rate. Alternatively, you can use an industry multiplier to multiply the business's EBITDA (Earnings Before Interest, Taxes, Depreciation, and Amortization) by a factor common in the industry. From the documents, it appears that the 2011 EBITDA multiplier for hospitality industries is 13.8. and in the end for small business in most sectors you can use a rule of thumb, 2x net + FFE.                                                                                                                    
\end{agentbubble}
\noindent\hfill{\footnotesize\textcolor{doccolor}{\textbf{$\hookrightarrow$\enspace Retrieved documents~(4)}}}
\begin{docenv}
{\footnotesize\textcolor{doccolor}{\textbf{Doc~1}}}\hfill{\tiny\texttt{score:~12.515}}\ {\tiny\textcolor{gray}{\{'relevant': \{'868c6e58-57c1-48ee-8c67-832dbedcac18': \{'value': 'no', 'timestamp': 1724063202\}\}\}}}\\[3pt]
{\small\textit{(sans titre)}}\\[2pt]
\vspace{2pt}
{\footnotesize 
You need experience. Work in retail, hopefully to some sort of managerial level where you get to see behind the scenes and not just manage associates.  If you're really gung-ho about it and willing to take a big risk, I've heard many people say the same reason that a small business fails - not enough cash flow. You might be lucky to raise enough capital for start-up costs, but can you survive two years of negative income? And I'm not even talking about your personal living expenses. Forget about that, you're going to be living at home eating ramen. I'm talking about having enough cash to pay the rent and keep the lights on. Unfortunately, it's going to be a vicious cycle because banks know small businesses are a huge risk and are unwilling to provide loans for them to survive...but small businesses fail because they aren't able to survive in the short-term no matter how good the prospects are in the long-term.   Your safest bet is to grow organically. This means having a business model with negligible start-up costs and little overhead costs. Your inventory is going to be small and because of that, your margins are going to be thin (no wholesale prices). But at this point, based on the products you're selling, you're basically just another one of the millions of internet re-sellers. You're going to basically have to place all your value in getting savvy to learn how to spot good deals and re-sell products for a higher price.   Honestly, I'll be harsh about it. You have 0 competitive advantage and offer 0 extra value in your idea. There's really no point in wasting your time and money unless you can pivot the idea so you're creating actual competitive advantages for yourself. Otherwise, just stick to re-selling online, at the very least to get a feel for the profit margins and gain experience.}
\end{docenv}
\begin{docenv}
{\footnotesize\textcolor{doccolor}{\textbf{Doc~2}}}\hfill{\tiny\texttt{score:~12.044}}\ {\tiny\textcolor{gray}{\{'relevant': \{'868c6e58-57c1-48ee-8c67-832dbedcac18': \{'value': 'yes', 'timestamp': 1724058177\}\}\}}}\\[3pt]
{\small\textit{(sans titre)}}\\[2pt]
\vspace{2pt}
{\footnotesize 
I don't have any experience in this, but this is my academic understanding of business pricing.  The LOWEST amount a seller would accept is the liquidation value. For a B\&B, what would the value of the land, the house, the furnishings, accounts payable, etc. be if it had to be sold today, minus any liabilities.  The amount the seller would like to pay for is going to be a multiple of its annual earnings. One example of this is the discounted cash flow analysis. You determine the EBITDA, the earnings a company generated, before interest,  depreciation, taxation and amortization.  Once you have this amount, you can project it out in perpetuity, or you use an industry multiplier.  Perpetuity:  You project this value out in perpituity, discounted by the going interest rate. In other words, if you project the business will earn \$100,000/year, the business should grow at a 5\% rate, and the going interest rate is 8\%. Using a growing perpetuity formula, one value of a business would be: 100,000 / (.08 - .03) = \$2,000,000.  This is a very high number, and the seller would love to get it.  It's more common to do a multiple of the EBIDTA. You can do some research into the valuation of the particular industry to figure out the EBIDTA multiplier for the industry. For example, this article suggests that the 2011 EBITDA multiplier for hospitality industries is 13.8. (It's valuing large hotel chains, but it's a start). So the value of this B\&B would be around \$1,380,000. Here is an online SME valuation tool to help with the EBIDTA multiple based valuation. Also, from my research, it looks like many small business use Seller Discretionary Earnings  (SDE) instead of EBITDA. I don't know much about it, but it seems to serve a similar purpose as EBITDA.  A potential buyer should request the financial statements of the business for the last few years to determine the value of the business, and then can negotiate with the owner a price. You would probably want to enlist a broker to help you with the transaction.}
\end{docenv}
\begin{docenv}
{\footnotesize\textcolor{doccolor}{\textbf{Doc~3}}}\hfill{\tiny\texttt{score:~11.899}}\ {\tiny\textcolor{gray}{\{'relevant': \{'868c6e58-57c1-48ee-8c67-832dbedcac18': \{'value': 'no', 'timestamp': 1724058210\}\}\}}}\\[3pt]
{\small\textit{(sans titre)}}\\[2pt]
\vspace{2pt}
{\footnotesize 
"I would suggest you forget everything you learned in economics. The only applicable knowledge is Accounting 101.  Step 1: An accrual basis financial statement. There is no step 2 if you don't do this. Most small business do everything cash basis. Simpler, cheaper but useless for analysis. You would get better answers from the local fortune teller than a cash basis statement.   Make one change from the general rules. If you have debt or are paying interest for inventory include that in your cost of sales. This is actually proper but the rule is little known and often ignored. Interest on debt up to the amount of inventory is a cost of inventory.  Step 2: Gross profit. If you seem to be working hard and still losing money it may be because you are selling products for less than they cost you. In this case the more you sell the more you lose. So suggestions like advertising or doing anything to increase sales are actually destructive.  Step 3 Price products at the level necessary to turn a profit at current sales and overhead. 'When we have enough sales we will make a profit"" is the philosophy of a start up business. It is toxic for a going concern.  Step 4 If sales are unsustainable at the price that produces a profit have the courage to sell or close the business. I have seen people waste their lives on futile endeavors just because they can't make that tough decision.  Finally Step 0: Ignore all other suggestions but this. They are well meaning but ill informed. To reiterate, growing sales while losing money on every transaction is a huge mistake. Trends, books, charts and graphs, analytics and market research are the tools of con-men and fortune tellers. Business is arithmetic and nothing more or less.  FYI if I don't get at least one upvote, this is the last time I am giving my valuable professional advice away for free on reddit. Folks will have to rely on the suggestions of their fellow college kids."}
\end{docenv}
\begin{docenv}
{\footnotesize\textcolor{doccolor}{\textbf{Doc~4}}}\hfill{\tiny\texttt{score:~13.988}}\ {\tiny\textcolor{gray}{\{'relevant': \{'868c6e58-57c1-48ee-8c67-832dbedcac18': \{'value': 'yes', 'timestamp': 1724058233\}\}\}}}\\[3pt]
{\small\textit{(sans titre)}}\\[2pt]
\vspace{2pt}
{\footnotesize 
nowhere near enough information to really help you.  Price it to an IRR of about 25\% in 4 years as your business has little to no history.  For small business in most sectors you can use a rule of thumb, 2x net + FFE.  For me personally, I wouldn't touch a business only 6 months old.}
\end{docenv}
\begin{refanswerenv}
{\footnotesize\textit{\textcolor{gray}{\textbf{Reference answer:}\enspace Based on the documents provided, there are a few approaches to pricing a small business. One approach is to determine the lowest amount a seller would accept, which is the liquidation value. This involves calculating the value of the business's assets, such as land, house, furnishings, and accounts payable, minus any liabilities. Another approach is to price the business based on a multiple of its annual earnings. This can be done using a discounted cash flow analysis, where you project the business's earnings out in perpetuity and discount it by the going interest rate. Alternatively, you can use an industry multiplier to multiply the business's EBITDA (Earnings Before Interest, Taxes, Depreciation, and Amortization) by a factor common in the industry. From the documents, it appears that the 2011 EBITDA multiplier for hospitality industries is 13.8. Additionally, many small businesses use Seller Discretionary Earnings (SDE) instead of EBITDA to determine the value of the business. A potential buyer should request the financial statements of the business for the last few years to determine the value of the business and then negotiate a price. It is also recommended to enlist a broker to help with the transaction.}}}
\end{refanswerenv}
\vspace{0.5em}
\noindent{\small\textcolor{usercolor}{\textbf{$\triangleright$\enspace User --- Turn~2}}}\hfill{\tiny\textcolor{gray}{Explanation $\cdot$ Follow-up $\cdot$ ANSWERABLE}}
\begin{userbubble}
What could be the reason that most small business fail along the way?
\end{userbubble}
\noindent\hfill{\small\textcolor{agentcolor}{\textbf{Agent\enspace$\triangleleft$}}}
\begin{agentbubble}
One of the reason most small businesses fail it is because banks usually know that small businesses are a huge risk and are unwilling to provide loans for them to survive. They aren't able to survive in the short-term no matter how good the prospects are in the long-term. They work crazy long hours and submit themselves to huge financial risk. Most owners have no interest in taking true responsibility for the company's health and in the end business has one imperative: to make money. Moreover some small owners business are lacked of some  intangible qualities of entrepreneurship. They  put in loads of commitment or loads of money into an idea that just isn't right. Or they spend money in the wrong ways.                                   
\end{agentbubble}
\noindent\hfill{\footnotesize\textcolor{doccolor}{\textbf{$\hookrightarrow$\enspace Retrieved documents~(7)}}}
\begin{docenv}
{\footnotesize\textcolor{doccolor}{\textbf{Doc~1}}}\hfill{\tiny\texttt{score:~53.667}}\ {\tiny\textcolor{gray}{\{'relevant': \{'868c6e58-57c1-48ee-8c67-832dbedcac18': \{'value': 'no', 'timestamp': 1724059149\}\}\}}}\\[3pt]
{\small\textit{(sans titre)}}\\[2pt]
\vspace{2pt}
{\footnotesize 
I don't have any experience in this, but this is my academic understanding of business pricing.  The LOWEST amount a seller would accept is the liquidation value. For a B\&B, what would the value of the land, the house, the furnishings, accounts payable, etc. be if it had to be sold today, minus any liabilities.  The amount the seller would like to pay for is going to be a multiple of its annual earnings. One example of this is the discounted cash flow analysis. You determine the EBITDA, the earnings a company generated, before interest,  depreciation, taxation and amortization.  Once you have this amount, you can project it out in perpetuity, or you use an industry multiplier.  Perpetuity:  You project this value out in perpituity, discounted by the going interest rate. In other words, if you project the business will earn \$100,000/year, the business should grow at a 5\% rate, and the going interest rate is 8\%. Using a growing perpetuity formula, one value of a business would be: 100,000 / (.08 - .03) = \$2,000,000.  This is a very high number, and the seller would love to get it.  It's more common to do a multiple of the EBIDTA. You can do some research into the valuation of the particular industry to figure out the EBIDTA multiplier for the industry. For example, this article suggests that the 2011 EBITDA multiplier for hospitality industries is 13.8. (It's valuing large hotel chains, but it's a start). So the value of this B\&B would be around \$1,380,000. Here is an online SME valuation tool to help with the EBIDTA multiple based valuation. Also, from my research, it looks like many small business use Seller Discretionary Earnings  (SDE) instead of EBITDA. I don't know much about it, but it seems to serve a similar purpose as EBITDA.  A potential buyer should request the financial statements of the business for the last few years to determine the value of the business, and then can negotiate with the owner a price. You would probably want to enlist a broker to help you with the transaction.}
\end{docenv}
\begin{docenv}
{\footnotesize\textcolor{doccolor}{\textbf{Doc~2}}}\hfill{\tiny\texttt{score:~28.674}}\ {\tiny\textcolor{gray}{\{'relevant': \{'868c6e58-57c1-48ee-8c67-832dbedcac18': \{'value': 'no', 'timestamp': 1724059136\}\}\}}}\\[3pt]
{\small\textit{(sans titre)}}\\[2pt]
\vspace{2pt}
{\footnotesize 
"EBITDA is in my opinion not a useful measure for an investor looking to buy shares on the stock market. It is more useful for private businesses open to changing their structuring, or looking to sell significant parts of their business. One of the main benefits of reporting Earnings Before Interest, Taxes, Depreciation \& Amortization, is that it presents the company as it would look to a potential buyer. Consider that net income, as a metric, includes interest costs, taxes, and depreciation. Interest costs are (to put it simply) a result of multiplying a business's debt by its interest rate. If you own a business, and personally guarantee the loan that the company has with the bank, your interest rates might be artificially low. If you have a policy of reaching high debt levels relative to your equity, in order to achieve high 'financial leveraging', your interest cost might be artificially high. Either way, if I bought your business, my debt structure could be completely different, and therefore your interest costs are not particularly relevant to me, a potential buyer. Instead, I should attempt to anticipate what my own interest costs would be, under my plans for your business. Taxes are a result of many factors, including the corporate structure of the business. If you run your business as a sole proprietorship (ie: no corporation), but I want to buy it under my corporation, then my tax rates could look nothing like yours. Or if we operated in multiple jurisdictions. etc. etc. Instead of using your taxes as an estimate for mine, I should anticipate my taxes based on my plans for your business. Depreciation / amortization is a measure that estimates how much of a business's ""fixed assets"" were ""used up"" during the year. ie: how much wear and tear occurred on your fleet of trucks? It is generally calculated as a \% of your overall asset value. It is a (very loose) proxy for the cash costs which will ultimately be incurred to make repairs/replacements. D\&A is also something which could significantly change if a business changes hands.}
\end{docenv}
\begin{docenv}
{\footnotesize\textcolor{doccolor}{\textbf{Doc~3}}}\hfill{\tiny\texttt{score:~27.412}}\ {\tiny\textcolor{gray}{\{'relevant': \{'868c6e58-57c1-48ee-8c67-832dbedcac18': \{'value': 'no', 'timestamp': 1724059126\}\}\}}}\\[3pt]
{\small\textit{(sans titre)}}\\[2pt]
\vspace{2pt}
{\footnotesize 
I'm not a finance guy by trade, I'm an accountant, so I'm not 100\% sure, but I'm going to say no.   EBITDA is your accounting profit before interest, taxes, depreciation and amortization. So it's already free of many non-cash items, and is closer to a cash-basis measure of profit, but still includes many non-cash revenues.   As far as I know, the discounting is to take the time value of money into account so you can make a decision on whether the present value of the expected future cash flows are greater than the present value of the cash outflow being asked today. If it is, then you buy. If not, then you don't.}
\end{docenv}
\begin{docenv}
{\footnotesize\textcolor{doccolor}{\textbf{Doc~4}}}\hfill{\tiny\texttt{score:~16.832}}\ {\tiny\textcolor{gray}{\{'relevant': \{'868c6e58-57c1-48ee-8c67-832dbedcac18': \{'value': 'no', 'timestamp': 1724059098\}\}\}}}\\[3pt]
{\small\textit{(sans titre)}}\\[2pt]
\vspace{2pt}
{\footnotesize 
\&gt;There have been a number of studies that have looked at this issue. This chart, from Web site designer smallbusinessplanned.com, summarizes the results of three different studies. Basically, after four years, 50 percent of the businesses are open. As time goes on, the success rate decreases, but it never gets to a failure rate of “nine out of 10.” \&gt;[...] \&gt;Even this does not show the whole picture. As Brian Headd, an economist at the Small Business Administration, demonstrated in a 2002 study for Small Business Economics, **about one–third of closed business were actually successful when they “failed.”** \&gt;“The significant proportion of businesses that closed while successful calls into question the use of ‘business closure’ as a meaningful measure of business outcome,” the study says. “It appears that **many owners may have executed a planned exit strategy,** closed a business without excess debt, sold a viable business, or retired from the work force.” Now that doesn't necessarily mean that Rand Paul's point is WRONG (he is chiefly talking about government investing in HIGHLY LEVERAGED, HIGH-RISK, HIGH-TECH businesses, which are a different story) -- but it does mean that the statistic he is citing (general business failure rate) is an urban-myth-falsehood, however commonly-believed, or commonly-restated."}
\end{docenv}
\begin{docenv}
{\footnotesize\textcolor{doccolor}{\textbf{Doc~5}}}\hfill{\tiny\texttt{score:~16.725}}\ {\tiny\textcolor{gray}{\{'relevant': \{'868c6e58-57c1-48ee-8c67-832dbedcac18': \{'value': 'yes', 'timestamp': 1724058773\}\}\}}}\\[3pt]
{\small\textit{(sans titre)}}\\[2pt]
\vspace{2pt}
{\footnotesize 
You need experience. Work in retail, hopefully to some sort of managerial level where you get to see behind the scenes and not just manage associates.  If you're really gung-ho about it and willing to take a big risk, I've heard many people say the same reason that a small business fails - not enough cash flow. You might be lucky to raise enough capital for start-up costs, but can you survive two years of negative income? And I'm not even talking about your personal living expenses. Forget about that, you're going to be living at home eating ramen. I'm talking about having enough cash to pay the rent and keep the lights on. Unfortunately, it's going to be a vicious cycle because banks know small businesses are a huge risk and are unwilling to provide loans for them to survive...but small businesses fail because they aren't able to survive in the short-term no matter how good the prospects are in the long-term.   Your safest bet is to grow organically. This means having a business model with negligible start-up costs and little overhead costs. Your inventory is going to be small and because of that, your margins are going to be thin (no wholesale prices). But at this point, based on the products you're selling, you're basically just another one of the millions of internet re-sellers. You're going to basically have to place all your value in getting savvy to learn how to spot good deals and re-sell products for a higher price.   Honestly, I'll be harsh about it. You have 0 competitive advantage and offer 0 extra value in your idea. There's really no point in wasting your time and money unless you can pivot the idea so you're creating actual competitive advantages for yourself. Otherwise, just stick to re-selling online, at the very least to get a feel for the profit margins and gain experience.}
\end{docenv}
\begin{docenv}
{\footnotesize\textcolor{doccolor}{\textbf{Doc~6}}}\hfill{\tiny\texttt{score:~16.307}}\ {\tiny\textcolor{gray}{\{'relevant': \{'868c6e58-57c1-48ee-8c67-832dbedcac18': \{'value': 'yes', 'timestamp': 1724058783\}\}\}}}\\[3pt]
{\small\textit{(sans titre)}}\\[2pt]
\vspace{2pt}
{\footnotesize 
People with failed businesses also often put in crazy long hours and submit themselves to huge financial risk. That's what you have to do just to start a business. It doesn't predict success.  Why do businesses fail?  I can now mostly predict if a business will fail because lots of people have ridiculous ideas or absolutely no way to make money (99\% of businesses at startup events will not be around in a year).  Luck doesn't predict the success of a business, but neither does hard work.  I started 3 businesses in the past 2 years and they are all making money.  I also put a ton of research into each one and carefully planned out a business strategy.    Many people starting businesses either don't have the ability to do this or just don't put the time into it because they are afraid that their awesome new idea won't be successful.}
\end{docenv}
\begin{docenv}
{\footnotesize\textcolor{doccolor}{\textbf{Doc~7}}}\hfill{\tiny\texttt{score:~13.818}}\ {\tiny\textcolor{gray}{\{'relevant': \{'868c6e58-57c1-48ee-8c67-832dbedcac18': \{'value': 'yes', 'timestamp': 1724058798\}\}\}}}\\[3pt]
{\small\textit{(sans titre)}}\\[2pt]
\vspace{2pt}
{\footnotesize 
"I'm sceptical of your position.   I reckon that people who put crazy long hours into a business that fails are lacking some intangible qualities of entrepreneurship.  There are people you can just tell, they don't have what it takes. They're putting in loads of commitment or loads of money into an idea that just isn't right. Or they spend money in the wrong ways - you know how all these internet companies spent money on champagne breakfasts and \$1000 chairs during the dot-com bubble? I've seen the small-scale version of that, as people rent expensive offices when I'm thinking ""why aren't you just working from home?""   As a (small time) entrepreneur I'm obviously biased, but if you've got what it takes (so I think there is an element of ""predestination"" here) and you work hard then you *will* succeed, to some greater or lesser extent.  I mean, I know electricians and builders that have started their own contracting or supplies businesses - they're not rich by the standard of mutli-multimillionaires, but they have a bloody nice house and a bloody nice car in the driveway and they probably take an afternoon a week off to play golf. They're doing better than most of the other guys they worked with as electricians or builders 20 years ago, and I don't see that they're really doing anything too clever or sophisticated.   Also, the classic figure of 9/10 businesses failed - ""failed"" here is a very dramatic word, implying that the entrepreneur lost his or her shirt. Often it's not a case that the business lost them money - they pissed away thousands on rent or wages or misplaced marketing - but more that it made a bit of money, but not enough, so they moved on to something else. That something else might be regular employment for some people, but there are *plenty* of entrepreneurs - perhaps some of the greatest and most famous ones - who have gone on from one failure to a much bigger success."}
\end{docenv}
\begin{refanswerenv}
{\footnotesize\textit{\textcolor{gray}{\textbf{Reference answer:}\enspace While the documents provided do not contain specific information about why most small businesses fail, there are several common reasons that small businesses often struggle. These can include lack of a solid business plan, insufficient capital, poor management, lack of market demand, and strong competition. It's important for small business owners to carefully consider these factors when starting and growing their businesses.}}}
\end{refanswerenv}
\vspace{0.5em}
\noindent{\small\textcolor{usercolor}{\textbf{$\triangleright$\enspace User --- Turn~3}}}\hfill{\tiny\textcolor{gray}{Opinion, Summarization $\cdot$ Follow-up $\cdot$ ANSWERABLE}}
\begin{userbubble}
What does it makes a business to go forward?
\end{userbubble}
\noindent\hfill{\small\textcolor{agentcolor}{\textbf{Agent\enspace$\triangleleft$}}}
\begin{agentbubble}
A business can move forward successfully by having a solid business model with minimal startup costs and low overhead costs. It's important to have a competitive advantage and offer extra value to customers. This can be achieved by becoming savvy in spotting good deals and reselling products for a higher price. Additionally you need to be careful with the structure, do all the research; understand the risk and opportunities as much you can.  Indeed to move forward in your business you need to make the things happen, you have to do the marketing and sale-link up with your prospective customers. And you have to have a well designed and well operating business. The key to good business management is ""do unto others as you would have them do unto you"" - ie. look at your business from your customers' perspective. In the end the point of a business is to make a profit . However, have some criteria for success or failure. Don't end up caught in the sunk costs fallacy. If the business isn't going to provide a better life than working, pull the plug early.
\end{agentbubble}
\noindent\hfill{\footnotesize\textcolor{doccolor}{\textbf{$\hookrightarrow$\enspace Retrieved documents~(6)}}}
\begin{docenv}
{\footnotesize\textcolor{doccolor}{\textbf{Doc~1}}}\hfill{\tiny\texttt{score:~48.075}}\ {\tiny\textcolor{gray}{\{'relevant': \{'868c6e58-57c1-48ee-8c67-832dbedcac18': \{'value': 'no', 'timestamp': 1724061208\}\}\}}}\\[3pt]
{\small\textit{(sans titre)}}\\[2pt]
\vspace{2pt}
{\footnotesize 
I don't have any experience in this, but this is my academic understanding of business pricing.  The LOWEST amount a seller would accept is the liquidation value. For a B\&B, what would the value of the land, the house, the furnishings, accounts payable, etc. be if it had to be sold today, minus any liabilities.  The amount the seller would like to pay for is going to be a multiple of its annual earnings. One example of this is the discounted cash flow analysis. You determine the EBITDA, the earnings a company generated, before interest,  depreciation, taxation and amortization.  Once you have this amount, you can project it out in perpetuity, or you use an industry multiplier.  Perpetuity:  You project this value out in perpituity, discounted by the going interest rate. In other words, if you project the business will earn \$100,000/year, the business should grow at a 5\% rate, and the going interest rate is 8\%. Using a growing perpetuity formula, one value of a business would be: 100,000 / (.08 - .03) = \$2,000,000.  This is a very high number, and the seller would love to get it.  It's more common to do a multiple of the EBIDTA. You can do some research into the valuation of the particular industry to figure out the EBIDTA multiplier for the industry. For example, this article suggests that the 2011 EBITDA multiplier for hospitality industries is 13.8. (It's valuing large hotel chains, but it's a start). So the value of this B\&B would be around \$1,380,000. Here is an online SME valuation tool to help with the EBIDTA multiple based valuation. Also, from my research, it looks like many small business use Seller Discretionary Earnings  (SDE) instead of EBITDA. I don't know much about it, but it seems to serve a similar purpose as EBITDA.  A potential buyer should request the financial statements of the business for the last few years to determine the value of the business, and then can negotiate with the owner a price. You would probably want to enlist a broker to help you with the transaction.}
\end{docenv}
\begin{docenv}
{\footnotesize\textcolor{doccolor}{\textbf{Doc~2}}}\hfill{\tiny\texttt{score:~27.560}}\ {\tiny\textcolor{gray}{\{'relevant': \{'868c6e58-57c1-48ee-8c67-832dbedcac18': \{'value': 'no', 'timestamp': 1724061211\}\}\}}}\\[3pt]
{\small\textit{(sans titre)}}\\[2pt]
\vspace{2pt}
{\footnotesize 
"EBITDA is in my opinion not a useful measure for an investor looking to buy shares on the stock market. It is more useful for private businesses open to changing their structuring, or looking to sell significant parts of their business. One of the main benefits of reporting Earnings Before Interest, Taxes, Depreciation \& Amortization, is that it presents the company as it would look to a potential buyer. Consider that net income, as a metric, includes interest costs, taxes, and depreciation. Interest costs are (to put it simply) a result of multiplying a business's debt by its interest rate. If you own a business, and personally guarantee the loan that the company has with the bank, your interest rates might be artificially low. If you have a policy of reaching high debt levels relative to your equity, in order to achieve high 'financial leveraging', your interest cost might be artificially high. Either way, if I bought your business, my debt structure could be completely different, and therefore your interest costs are not particularly relevant to me, a potential buyer. Instead, I should attempt to anticipate what my own interest costs would be, under my plans for your business. Taxes are a result of many factors, including the corporate structure of the business. If you run your business as a sole proprietorship (ie: no corporation), but I want to buy it under my corporation, then my tax rates could look nothing like yours. Or if we operated in multiple jurisdictions. etc. etc. Instead of using your taxes as an estimate for mine, I should anticipate my taxes based on my plans for your business. Depreciation / amortization is a measure that estimates how much of a business's ""fixed assets"" were ""used up"" during the year. ie: how much wear and tear occurred on your fleet of trucks? It is generally calculated as a \% of your overall asset value. It is a (very loose) proxy for the cash costs which will ultimately be incurred to make repairs/replacements. D\&A is also something which could significantly change if a business changes hands.}
\end{docenv}
\begin{docenv}
{\footnotesize\textcolor{doccolor}{\textbf{Doc~3}}}\hfill{\tiny\texttt{score:~27.556}}\ {\tiny\textcolor{gray}{\{'relevant': \{'868c6e58-57c1-48ee-8c67-832dbedcac18': \{'value': 'yes', 'timestamp': 1724061221\}\}\}}}\\[3pt]
{\small\textit{(sans titre)}}\\[2pt]
\vspace{2pt}
{\footnotesize 
You need experience. Work in retail, hopefully to some sort of managerial level where you get to see behind the scenes and not just manage associates.  If you're really gung-ho about it and willing to take a big risk, I've heard many people say the same reason that a small business fails - not enough cash flow. You might be lucky to raise enough capital for start-up costs, but can you survive two years of negative income? And I'm not even talking about your personal living expenses. Forget about that, you're going to be living at home eating ramen. I'm talking about having enough cash to pay the rent and keep the lights on. Unfortunately, it's going to be a vicious cycle because banks know small businesses are a huge risk and are unwilling to provide loans for them to survive...but small businesses fail because they aren't able to survive in the short-term no matter how good the prospects are in the long-term.   Your safest bet is to grow organically. This means having a business model with negligible start-up costs and little overhead costs. Your inventory is going to be small and because of that, your margins are going to be thin (no wholesale prices). But at this point, based on the products you're selling, you're basically just another one of the millions of internet re-sellers. You're going to basically have to place all your value in getting savvy to learn how to spot good deals and re-sell products for a higher price.   Honestly, I'll be harsh about it. You have 0 competitive advantage and offer 0 extra value in your idea. There's really no point in wasting your time and money unless you can pivot the idea so you're creating actual competitive advantages for yourself. Otherwise, just stick to re-selling online, at the very least to get a feel for the profit margins and gain experience.}
\end{docenv}
\begin{docenv}
{\footnotesize\textcolor{doccolor}{\textbf{Doc~4}}}\hfill{\tiny\texttt{score:~13.046}}\ {\tiny\textcolor{gray}{\{'relevant': \{'868c6e58-57c1-48ee-8c67-832dbedcac18': \{'value': 'yes', 'timestamp': 1724061967\}\}\}}}\\[3pt]
{\small\textit{(sans titre)}}\\[2pt]
\vspace{2pt}
{\footnotesize 
Going into a business with the potential to grow is the high risk/high reward option. If it pans out, you could do brilliantly, but there's a very real chance you'll work like a dog for years and have nothing to show for it at the end other than debt. Even if the business is successful, if you aren't careful with the structure, you might not benefit.  Do all the research you can, so you can understand the risks and opportunities as much as possible. Then you can evaluate whether it's worth it.  Welding is something you can always go back to if the business doesn't pan out. If you go with the business plan, have some criteria for success or failure. Don't end up caught in the sunk costs fallacy. If the business isn't going to provide a better life than working, pull the plug early.}
\end{docenv}
\begin{docenv}
{\footnotesize\textcolor{doccolor}{\textbf{Doc~5}}}\hfill{\tiny\texttt{score:~12.767}}\ {\tiny\textcolor{gray}{\{'relevant': \{'868c6e58-57c1-48ee-8c67-832dbedcac18': \{'value': 'yes', 'timestamp': 1724062683\}\}\}}}\\[3pt]
{\small\textit{(sans titre)}}\\[2pt]
\vspace{2pt}
{\footnotesize 
"Things don't just happen, you have to make them happen, you have to do the marketing and sales - link up with your prospective customers. And you have to have a well designed and well operating business. The key to good business management is ""do unto others as you would have them do unto you"" - ie. look at your business from your customers' perspective."}
\end{docenv}
\begin{docenv}
{\footnotesize\textcolor{doccolor}{\textbf{Doc~6}}}\hfill{\tiny\texttt{score:~11.702}}\ {\tiny\textcolor{gray}{\{'relevant': \{'868c6e58-57c1-48ee-8c67-832dbedcac18': \{'value': 'yes', 'timestamp': 1724062750\}\}\}}}\\[3pt]
{\small\textit{(sans titre)}}\\[2pt]
\vspace{2pt}
{\footnotesize 
\&gt;The point of a business is to make a profit. End of story.  Thats a narrow and short-sighted interpretation of business practices. The rest of your comment devolved into strawmen and bullshit so I'm not going to bother addressing it. Your links were equally irrelevant but I appreciate the fact that you at least *tried* to source a reasoned argument. You didn't succeed but at least you tried.}
\end{docenv}
\begin{refanswerenv}
{\footnotesize\textit{\textcolor{gray}{\textbf{Reference answer:}\enspace A business can move forward successfully by having a solid business model with minimal startup costs and low overhead costs. It's important to have a competitive advantage and offer extra value to customers. This can be achieved by becoming savvy in spotting good deals and reselling products for a higher price. Additionally, having a clear understanding of the financials, such as calculating the liquidation value or using industry multipliers for pricing, can contribute to the business's success. Furthermore, entrepreneurs should be willing to take responsibility for the company's health, commit to the right ideas, and spend money wisely.}}}
\end{refanswerenv}
\vspace{0.5em}
\noindent{\small\textcolor{usercolor}{\textbf{$\triangleright$\enspace User --- Turn~4}}}\hfill{\tiny\textcolor{gray}{Opinion, Summarization $\cdot$ Follow-up $\cdot$ ANSWERABLE}}
\begin{userbubble}
Speaking about good and bad business, what does it make a good manager?
\end{userbubble}
\noindent\hfill{\small\textcolor{agentcolor}{\textbf{Agent\enspace$\triangleleft$}}}
\begin{agentbubble}
One of the first things that a good manager must have is empathy. The last thing you want is unhappy employees. You can gain their confidence by being good at your job and making their job feel better. A good manager can find out what people's real skills are then design the role around them. They don't force people into strictly defined roles. It is essential to learn how to do the job and don't be afraid to step in and help out. If you never do this people will assume you think you're too good to do the work they do and they won't respect you or your authority. Being good at something gives you natural authority and people will seek you out for advice/direction. Trust your employees to do their jobs until they prove otherwise. Management is a relationship and relationships are based on trust. A manager should be much more aware of how each person contributes to the whole, what their strengths are, who to talk to about what, and how, and why. Moreover  good managers have the most immediate influence over the comfort of their employees while on the job. Comfortable office furniture, equipment that doesn't suck, amenities to make being at work more of a pleasure - all means of showing employees they're treated as valuable professionals.
\end{agentbubble}
\noindent\hfill{\footnotesize\textcolor{doccolor}{\textbf{$\hookrightarrow$\enspace Retrieved documents~(7)}}}
\begin{docenv}
{\footnotesize\textcolor{doccolor}{\textbf{Doc~1}}}\hfill{\tiny\texttt{score:~42.686}}\ {\tiny\textcolor{gray}{\{'relevant': \{'868c6e58-57c1-48ee-8c67-832dbedcac18': \{'value': 'no', 'timestamp': 1724065449\}\}\}}}\\[3pt]
{\small\textit{(sans titre)}}\\[2pt]
\vspace{2pt}
{\footnotesize 
I don't have any experience in this, but this is my academic understanding of business pricing.  The LOWEST amount a seller would accept is the liquidation value. For a B\&B, what would the value of the land, the house, the furnishings, accounts payable, etc. be if it had to be sold today, minus any liabilities.  The amount the seller would like to pay for is going to be a multiple of its annual earnings. One example of this is the discounted cash flow analysis. You determine the EBITDA, the earnings a company generated, before interest,  depreciation, taxation and amortization.  Once you have this amount, you can project it out in perpetuity, or you use an industry multiplier.  Perpetuity:  You project this value out in perpituity, discounted by the going interest rate. In other words, if you project the business will earn \$100,000/year, the business should grow at a 5\% rate, and the going interest rate is 8\%. Using a growing perpetuity formula, one value of a business would be: 100,000 / (.08 - .03) = \$2,000,000.  This is a very high number, and the seller would love to get it.  It's more common to do a multiple of the EBIDTA. You can do some research into the valuation of the particular industry to figure out the EBIDTA multiplier for the industry. For example, this article suggests that the 2011 EBITDA multiplier for hospitality industries is 13.8. (It's valuing large hotel chains, but it's a start). So the value of this B\&B would be around \$1,380,000. Here is an online SME valuation tool to help with the EBIDTA multiple based valuation. Also, from my research, it looks like many small business use Seller Discretionary Earnings  (SDE) instead of EBITDA. I don't know much about it, but it seems to serve a similar purpose as EBITDA.  A potential buyer should request the financial statements of the business for the last few years to determine the value of the business, and then can negotiate with the owner a price. You would probably want to enlist a broker to help you with the transaction.}
\end{docenv}
\begin{docenv}
{\footnotesize\textcolor{doccolor}{\textbf{Doc~2}}}\hfill{\tiny\texttt{score:~30.902}}\ {\tiny\textcolor{gray}{\{'relevant': \{'868c6e58-57c1-48ee-8c67-832dbedcac18': \{'value': 'no', 'timestamp': 1724065452\}\}\}}}\\[3pt]
{\small\textit{(sans titre)}}\\[2pt]
\vspace{2pt}
{\footnotesize 
You need experience. Work in retail, hopefully to some sort of managerial level where you get to see behind the scenes and not just manage associates.  If you're really gung-ho about it and willing to take a big risk, I've heard many people say the same reason that a small business fails - not enough cash flow. You might be lucky to raise enough capital for start-up costs, but can you survive two years of negative income? And I'm not even talking about your personal living expenses. Forget about that, you're going to be living at home eating ramen. I'm talking about having enough cash to pay the rent and keep the lights on. Unfortunately, it's going to be a vicious cycle because banks know small businesses are a huge risk and are unwilling to provide loans for them to survive...but small businesses fail because they aren't able to survive in the short-term no matter how good the prospects are in the long-term.   Your safest bet is to grow organically. This means having a business model with negligible start-up costs and little overhead costs. Your inventory is going to be small and because of that, your margins are going to be thin (no wholesale prices). But at this point, based on the products you're selling, you're basically just another one of the millions of internet re-sellers. You're going to basically have to place all your value in getting savvy to learn how to spot good deals and re-sell products for a higher price.   Honestly, I'll be harsh about it. You have 0 competitive advantage and offer 0 extra value in your idea. There's really no point in wasting your time and money unless you can pivot the idea so you're creating actual competitive advantages for yourself. Otherwise, just stick to re-selling online, at the very least to get a feel for the profit margins and gain experience.}
\end{docenv}
\begin{docenv}
{\footnotesize\textcolor{doccolor}{\textbf{Doc~3}}}\hfill{\tiny\texttt{score:~26.089}}\ {\tiny\textcolor{gray}{\{'relevant': \{'868c6e58-57c1-48ee-8c67-832dbedcac18': \{'value': 'no', 'timestamp': 1724065456\}\}\}}}\\[3pt]
{\small\textit{(sans titre)}}\\[2pt]
\vspace{2pt}
{\footnotesize 
"EBITDA is in my opinion not a useful measure for an investor looking to buy shares on the stock market. It is more useful for private businesses open to changing their structuring, or looking to sell significant parts of their business. One of the main benefits of reporting Earnings Before Interest, Taxes, Depreciation \& Amortization, is that it presents the company as it would look to a potential buyer. Consider that net income, as a metric, includes interest costs, taxes, and depreciation. Interest costs are (to put it simply) a result of multiplying a business's debt by its interest rate. If you own a business, and personally guarantee the loan that the company has with the bank, your interest rates might be artificially low. If you have a policy of reaching high debt levels relative to your equity, in order to achieve high 'financial leveraging', your interest cost might be artificially high. Either way, if I bought your business, my debt structure could be completely different, and therefore your interest costs are not particularly relevant to me, a potential buyer. Instead, I should attempt to anticipate what my own interest costs would be, under my plans for your business. Taxes are a result of many factors, including the corporate structure of the business. If you run your business as a sole proprietorship (ie: no corporation), but I want to buy it under my corporation, then my tax rates could look nothing like yours. Or if we operated in multiple jurisdictions. etc. etc. Instead of using your taxes as an estimate for mine, I should anticipate my taxes based on my plans for your business. Depreciation / amortization is a measure that estimates how much of a business's ""fixed assets"" were ""used up"" during the year. ie: how much wear and tear occurred on your fleet of trucks? It is generally calculated as a \% of your overall asset value. It is a (very loose) proxy for the cash costs which will ultimately be incurred to make repairs/replacements. D\&A is also something which could significantly change if a business changes hands.}
\end{docenv}
\begin{docenv}
{\footnotesize\textcolor{doccolor}{\textbf{Doc~4}}}\hfill{\tiny\texttt{score:~17.794}}\ {\tiny\textcolor{gray}{\{'relevant': \{'868c6e58-57c1-48ee-8c67-832dbedcac18': \{'value': 'yes', 'timestamp': 1724065794\}\}\}}}\\[3pt]
{\small\textit{(sans titre)}}\\[2pt]
\vspace{2pt}
{\footnotesize 
A good manager sets up systems, so you hardly have do to direct anything yourself. Manage by exception. Trust,  ask employees how things can be improved(they will have better knowledge of the task) and let them implement it if it's good.  At our place,  we have visual management things on TV screens so everyone knows what's going on,  metrics etc.  People are constantly trying to improve the performance of system without management pushing for it.   Empathy is highly important in management.  You have to feel what they feel.  The last thing you want is disgruntled employees.  Good luck getting people to do things well,  when they think you are a dick.   You can't demand respect... By definition.  You gradually gain it by being good at the job,  and making their jobs feel better.  If you can find out what people's real skills are then design the role around them. Don't force people into strictly defined roles. This is why people get fired. They're in a job that doesn't match their skills. Nothing worse then knowing your shit at a job,  but unable to move to a job you think you are good at.}
\end{docenv}
\begin{docenv}
{\footnotesize\textcolor{doccolor}{\textbf{Doc~5}}}\hfill{\tiny\texttt{score:~15.665}}\ {\tiny\textcolor{gray}{\{'relevant': \{'868c6e58-57c1-48ee-8c67-832dbedcac18': \{'value': 'yes', 'timestamp': 1724065874\}\}\}}}\\[3pt]
{\small\textit{(sans titre)}}\\[2pt]
\vspace{2pt}
{\footnotesize 
Learn how to do the job and don't be afraid to step in and help out. If you never do this people will assume you think you're too good to do the work they do and they won't respect you or your authority.   Being good at something gives you natural authority and people will seek you out for advice/direction.  Trust your employees to do their jobs until they prove otherwise. Management is a relationship and relationships are based on trust. If you want to be untrusted the very best way to go about it is to show you don't trust the other person. If you're not trusted then no one is going to follow your untrustworthy advice/direction.  Stay calm even when things are going terribly. Even if you're freaking out on the inside it's best to show a calm front. If you're constantly on edge it will amplify through your employees and a chaotic workplace is obviously less efficient overtime. If you lose your cool apologize sincerely.  Be positive. Always have something good to say about something and employees will more likely be able to stay positive. If you're negative all the time then again this is going to amplify through employees and if everyone is gossiping and complaining all day efficiency is going to go down.   Be assertive and explain why. Sometimes you can take input from employees and negotiate, but sometimes you just have to tell them what to do and it's important to include why because if you include the why people are much more likely to comply as well as to remember what you said.  All these things may not make you liked. If your goal is to be liked then quit because all bosses aren't liked by many employees. You will be more likely to be effective and at least respected.}
\end{docenv}
\begin{docenv}
{\footnotesize\textcolor{doccolor}{\textbf{Doc~6}}}\hfill{\tiny\texttt{score:~14.229}}\ {\tiny\textcolor{gray}{\{'relevant': \{'868c6e58-57c1-48ee-8c67-832dbedcac18': \{'value': 'yes', 'timestamp': 1724065909\}\}\}}}\\[3pt]
{\small\textit{(sans titre)}}\\[2pt]
\vspace{2pt}
{\footnotesize 
"Theres a difference between being ""materially better"" and having deep knowledge of company politics as part of your job.   A manager should be much more aware of how each person contributes to the whole, what their stengths are, who to talk to about what, and how, and why. Thats why non-managerial level people are called ""individual contributors"" -- they're focused on doing their part. Where-as managers are focused on how all of the individual contributors work together.  Its not about being ""materially better"" it's about delegation of task types and workflows."}
\end{docenv}
\begin{docenv}
{\footnotesize\textcolor{doccolor}{\textbf{Doc~7}}}\hfill{\tiny\texttt{score:~14.081}}\ {\tiny\textcolor{gray}{\{'relevant': \{'868c6e58-57c1-48ee-8c67-832dbedcac18': \{'value': 'yes', 'timestamp': 1724066058\}\}\}}}\\[3pt]
{\small\textit{(sans titre)}}\\[2pt]
\vspace{2pt}
{\footnotesize 
There's no downplaying the importance of something as simple as that.  But I think that Future covers it:  \&gt; Employees care about your business… but only when you prove you care about them first.  Managers have the most immediate influence over the comfort of their employees while on the job. Comfortable office furniture, equipment that doesn't suck, amenities to make being at work more of a pleasure - all means of showing employees they're treated as valuable professionals.}
\end{docenv}
\begin{refanswerenv}
{\footnotesize\textit{\textcolor{gray}{\textbf{Reference answer:}\enspace A good manager is someone who leads their team effectively, communicates clearly, and makes informed decisions. They should be able to motivate their team, delegate tasks appropriately, and provide constructive feedback. A good manager also understands the importance of work-life balance and fosters a positive work environment. They should have a strong understanding of the business's goals and strategies, and be able to align their team's efforts with these objectives. Additionally, a good manager is open to new ideas, adaptable to change, and committed to continuous learning and improvement.}}}
\end{refanswerenv}
\vspace{0.5em}
\noindent{\small\textcolor{usercolor}{\textbf{$\triangleright$\enspace User --- Turn~5}}}\hfill{\tiny\textcolor{gray}{How-To $\cdot$ Follow-up $\cdot$ ANSWERABLE}}
\begin{userbubble}
How do companies hire their employees?
\end{userbubble}
\noindent\hfill{\small\textcolor{agentcolor}{\textbf{Agent\enspace$\triangleleft$}}}
\begin{agentbubble}
   It depends what kind the company  is looking for hire someone. Sometimes the school that the person attended is a good start. Many companies are looking for skills as marketing, leasing space, setting up a business, and management . Also a good attitude, culture can help . Unfortunately the problem is that most recruiters don't know what they are looking for either. They don't have the background or experience to match someone to a job. Or they ignore transferable skills. 
\end{agentbubble}
\noindent\hfill{\footnotesize\textcolor{doccolor}{\textbf{$\hookrightarrow$\enspace Retrieved documents~(6)}}}
\begin{docenv}
{\footnotesize\textcolor{doccolor}{\textbf{Doc~1}}}\hfill{\tiny\texttt{score:~42.686}}\ {\tiny\textcolor{gray}{\{'relevant': \{'868c6e58-57c1-48ee-8c67-832dbedcac18': \{'value': 'no', 'timestamp': 1724069832\}\}\}}}\\[3pt]
{\small\textit{(sans titre)}}\\[2pt]
\vspace{2pt}
{\footnotesize 
I don't have any experience in this, but this is my academic understanding of business pricing.  The LOWEST amount a seller would accept is the liquidation value. For a B\&B, what would the value of the land, the house, the furnishings, accounts payable, etc. be if it had to be sold today, minus any liabilities.  The amount the seller would like to pay for is going to be a multiple of its annual earnings. One example of this is the discounted cash flow analysis. You determine the EBITDA, the earnings a company generated, before interest,  depreciation, taxation and amortization.  Once you have this amount, you can project it out in perpetuity, or you use an industry multiplier.  Perpetuity:  You project this value out in perpituity, discounted by the going interest rate. In other words, if you project the business will earn \$100,000/year, the business should grow at a 5\% rate, and the going interest rate is 8\%. Using a growing perpetuity formula, one value of a business would be: 100,000 / (.08 - .03) = \$2,000,000.  This is a very high number, and the seller would love to get it.  It's more common to do a multiple of the EBIDTA. You can do some research into the valuation of the particular industry to figure out the EBIDTA multiplier for the industry. For example, this article suggests that the 2011 EBITDA multiplier for hospitality industries is 13.8. (It's valuing large hotel chains, but it's a start). So the value of this B\&B would be around \$1,380,000. Here is an online SME valuation tool to help with the EBIDTA multiple based valuation. Also, from my research, it looks like many small business use Seller Discretionary Earnings  (SDE) instead of EBITDA. I don't know much about it, but it seems to serve a similar purpose as EBITDA.  A potential buyer should request the financial statements of the business for the last few years to determine the value of the business, and then can negotiate with the owner a price. You would probably want to enlist a broker to help you with the transaction.}
\end{docenv}
\begin{docenv}
{\footnotesize\textcolor{doccolor}{\textbf{Doc~2}}}\hfill{\tiny\texttt{score:~30.902}}\ {\tiny\textcolor{gray}{\{'relevant': \{'868c6e58-57c1-48ee-8c67-832dbedcac18': \{'value': 'no', 'timestamp': 1724069834\}\}\}}}\\[3pt]
{\small\textit{(sans titre)}}\\[2pt]
\vspace{2pt}
{\footnotesize 
You need experience. Work in retail, hopefully to some sort of managerial level where you get to see behind the scenes and not just manage associates.  If you're really gung-ho about it and willing to take a big risk, I've heard many people say the same reason that a small business fails - not enough cash flow. You might be lucky to raise enough capital for start-up costs, but can you survive two years of negative income? And I'm not even talking about your personal living expenses. Forget about that, you're going to be living at home eating ramen. I'm talking about having enough cash to pay the rent and keep the lights on. Unfortunately, it's going to be a vicious cycle because banks know small businesses are a huge risk and are unwilling to provide loans for them to survive...but small businesses fail because they aren't able to survive in the short-term no matter how good the prospects are in the long-term.   Your safest bet is to grow organically. This means having a business model with negligible start-up costs and little overhead costs. Your inventory is going to be small and because of that, your margins are going to be thin (no wholesale prices). But at this point, based on the products you're selling, you're basically just another one of the millions of internet re-sellers. You're going to basically have to place all your value in getting savvy to learn how to spot good deals and re-sell products for a higher price.   Honestly, I'll be harsh about it. You have 0 competitive advantage and offer 0 extra value in your idea. There's really no point in wasting your time and money unless you can pivot the idea so you're creating actual competitive advantages for yourself. Otherwise, just stick to re-selling online, at the very least to get a feel for the profit margins and gain experience.}
\end{docenv}
\begin{docenv}
{\footnotesize\textcolor{doccolor}{\textbf{Doc~3}}}\hfill{\tiny\texttt{score:~26.089}}\ {\tiny\textcolor{gray}{\{'relevant': \{'868c6e58-57c1-48ee-8c67-832dbedcac18': \{'value': 'no', 'timestamp': 1724069836\}\}\}}}\\[3pt]
{\small\textit{(sans titre)}}\\[2pt]
\vspace{2pt}
{\footnotesize 
"EBITDA is in my opinion not a useful measure for an investor looking to buy shares on the stock market. It is more useful for private businesses open to changing their structuring, or looking to sell significant parts of their business. One of the main benefits of reporting Earnings Before Interest, Taxes, Depreciation \& Amortization, is that it presents the company as it would look to a potential buyer. Consider that net income, as a metric, includes interest costs, taxes, and depreciation. Interest costs are (to put it simply) a result of multiplying a business's debt by its interest rate. If you own a business, and personally guarantee the loan that the company has with the bank, your interest rates might be artificially low. If you have a policy of reaching high debt levels relative to your equity, in order to achieve high 'financial leveraging', your interest cost might be artificially high. Either way, if I bought your business, my debt structure could be completely different, and therefore your interest costs are not particularly relevant to me, a potential buyer. Instead, I should attempt to anticipate what my own interest costs would be, under my plans for your business. Taxes are a result of many factors, including the corporate structure of the business. If you run your business as a sole proprietorship (ie: no corporation), but I want to buy it under my corporation, then my tax rates could look nothing like yours. Or if we operated in multiple jurisdictions. etc. etc. Instead of using your taxes as an estimate for mine, I should anticipate my taxes based on my plans for your business. Depreciation / amortization is a measure that estimates how much of a business's ""fixed assets"" were ""used up"" during the year. ie: how much wear and tear occurred on your fleet of trucks? It is generally calculated as a \% of your overall asset value. It is a (very loose) proxy for the cash costs which will ultimately be incurred to make repairs/replacements. D\&A is also something which could significantly change if a business changes hands.}
\end{docenv}
\begin{docenv}
{\footnotesize\textcolor{doccolor}{\textbf{Doc~4}}}\hfill{\tiny\texttt{score:~18.429}}\ {\tiny\textcolor{gray}{\{'relevant': \{'868c6e58-57c1-48ee-8c67-832dbedcac18': \{'value': 'yes', 'timestamp': 1724070045\}\}\}}}\\[3pt]
{\small\textit{(sans titre)}}\\[2pt]
\vspace{2pt}
{\footnotesize 
What's easy? Debt, personal, marketing related finance.   What we look for in hires is simple though - modeling skills (everyone thinks they're better than they actually are) and case analysis (same story). If you can get those two under your belt, and a good school under your name you should be good.}
\end{docenv}
\begin{docenv}
{\footnotesize\textcolor{doccolor}{\textbf{Doc~5}}}\hfill{\tiny\texttt{score:~15.613}}\ {\tiny\textcolor{gray}{\{'relevant': \{'868c6e58-57c1-48ee-8c67-832dbedcac18': \{'value': 'yes', 'timestamp': 1724070190\}\}\}}}\\[3pt]
{\small\textit{(sans titre)}}\\[2pt]
\vspace{2pt}
{\footnotesize 
Unfortunately, I'm not the guy that makes the call about hires.  You should talk to the institutional guys - get to know them, find out where they go for drinks and show up and just shoot the shit and chit chat. Be outgoing and joke around.  Mention you'd like a shot and to keep you in mind if something opens up. I can tell you that in this industry the first thing people look for is if the candidate is a good fit for the team - both in terms of attitude, culture and overall fit.   Leverage what you have.  They already know you.}
\end{docenv}
\begin{docenv}
{\footnotesize\textcolor{doccolor}{\textbf{Doc~6}}}\hfill{\tiny\texttt{score:~15.831}}\ {\tiny\textcolor{gray}{\{'relevant': \{'868c6e58-57c1-48ee-8c67-832dbedcac18': \{'value': 'yes', 'timestamp': 1724070233\}\}\}}}\\[3pt]
{\small\textit{(sans titre)}}\\[2pt]
\vspace{2pt}
{\footnotesize 
The problem is that most recruiters don't know what they are looking for either. They don't have the background or experience to match someone to a job. Or they ignore transferable skills. The talent acquisition department at my company had such a bad reputation everyone says to just email the hiring manager.}
\end{docenv}
\begin{refanswerenv}
{\footnotesize\textit{\textcolor{gray}{\textbf{Reference answer:}\enspace Companies typically hire employees through a process that involves posting job openings, reviewing resumes and applications, conducting interviews, and checking references. This process can vary depending on the size and type of the company, as well as the specific role being filled. Some companies may also use recruiting agencies or online platforms to help find and vet potential candidates. Once a candidate is selected, they typically go through an onboarding process to learn about the company's culture, policies, and procedures.}}}
\end{refanswerenv}
\vspace{0.5em}
\noindent{\small\textcolor{usercolor}{\textbf{$\triangleright$\enspace User --- Turn~6}}}\hfill{\tiny\textcolor{gray}{Factoid, Summarization $\cdot$ Follow-up $\cdot$ ANSWERABLE}}
\begin{userbubble}
What are the basic benefits that a company to do pay its employee besides salary, of course?
\end{userbubble}
\noindent\hfill{\small\textcolor{agentcolor}{\textbf{Agent\enspace$\triangleleft$}}}
\begin{agentbubble}
Actually  companies are not required to offer their employees any compensation above paying your salary. However, benefits attract higher quality employees.  Employees want certain jobs because of the tax benefits. Employees expect  the benefits because almost all reputable companies of any significant size offer it. You could run a great company, but if you don't offer a 401k plan, you can scare away good potential employees. It would give a bad impression the same way that not offering health insurance would. Some companies offer better benefits than others. For instance,  Apple offers 15\% discount on stock, 25\% 401K matching and plus excellent health/vision/dental/etc. benefit. Additionally personnel training  is a great benefit also to the employees.    
\end{agentbubble}
\noindent\hfill{\footnotesize\textcolor{doccolor}{\textbf{$\hookrightarrow$\enspace Retrieved documents~(6)}}}
\begin{docenv}
{\footnotesize\textcolor{doccolor}{\textbf{Doc~1}}}\hfill{\tiny\texttt{score:~42.686}}\ {\tiny\textcolor{gray}{\{'relevant': \{'868c6e58-57c1-48ee-8c67-832dbedcac18': \{'value': 'no', 'timestamp': 1724071414\}\}\}}}\\[3pt]
{\small\textit{(sans titre)}}\\[2pt]
\vspace{2pt}
{\footnotesize 
I don't have any experience in this, but this is my academic understanding of business pricing.  The LOWEST amount a seller would accept is the liquidation value. For a B\&B, what would the value of the land, the house, the furnishings, accounts payable, etc. be if it had to be sold today, minus any liabilities.  The amount the seller would like to pay for is going to be a multiple of its annual earnings. One example of this is the discounted cash flow analysis. You determine the EBITDA, the earnings a company generated, before interest,  depreciation, taxation and amortization.  Once you have this amount, you can project it out in perpetuity, or you use an industry multiplier.  Perpetuity:  You project this value out in perpituity, discounted by the going interest rate. In other words, if you project the business will earn \$100,000/year, the business should grow at a 5\% rate, and the going interest rate is 8\%. Using a growing perpetuity formula, one value of a business would be: 100,000 / (.08 - .03) = \$2,000,000.  This is a very high number, and the seller would love to get it.  It's more common to do a multiple of the EBIDTA. You can do some research into the valuation of the particular industry to figure out the EBIDTA multiplier for the industry. For example, this article suggests that the 2011 EBITDA multiplier for hospitality industries is 13.8. (It's valuing large hotel chains, but it's a start). So the value of this B\&B would be around \$1,380,000. Here is an online SME valuation tool to help with the EBIDTA multiple based valuation. Also, from my research, it looks like many small business use Seller Discretionary Earnings  (SDE) instead of EBITDA. I don't know much about it, but it seems to serve a similar purpose as EBITDA.  A potential buyer should request the financial statements of the business for the last few years to determine the value of the business, and then can negotiate with the owner a price. You would probably want to enlist a broker to help you with the transaction.}
\end{docenv}
\begin{docenv}
{\footnotesize\textcolor{doccolor}{\textbf{Doc~2}}}\hfill{\tiny\texttt{score:~30.902}}\ {\tiny\textcolor{gray}{\{'relevant': \{'868c6e58-57c1-48ee-8c67-832dbedcac18': \{'value': 'no', 'timestamp': 1724071416\}\}\}}}\\[3pt]
{\small\textit{(sans titre)}}\\[2pt]
\vspace{2pt}
{\footnotesize 
You need experience. Work in retail, hopefully to some sort of managerial level where you get to see behind the scenes and not just manage associates.  If you're really gung-ho about it and willing to take a big risk, I've heard many people say the same reason that a small business fails - not enough cash flow. You might be lucky to raise enough capital for start-up costs, but can you survive two years of negative income? And I'm not even talking about your personal living expenses. Forget about that, you're going to be living at home eating ramen. I'm talking about having enough cash to pay the rent and keep the lights on. Unfortunately, it's going to be a vicious cycle because banks know small businesses are a huge risk and are unwilling to provide loans for them to survive...but small businesses fail because they aren't able to survive in the short-term no matter how good the prospects are in the long-term.   Your safest bet is to grow organically. This means having a business model with negligible start-up costs and little overhead costs. Your inventory is going to be small and because of that, your margins are going to be thin (no wholesale prices). But at this point, based on the products you're selling, you're basically just another one of the millions of internet re-sellers. You're going to basically have to place all your value in getting savvy to learn how to spot good deals and re-sell products for a higher price.   Honestly, I'll be harsh about it. You have 0 competitive advantage and offer 0 extra value in your idea. There's really no point in wasting your time and money unless you can pivot the idea so you're creating actual competitive advantages for yourself. Otherwise, just stick to re-selling online, at the very least to get a feel for the profit margins and gain experience.}
\end{docenv}
\begin{docenv}
{\footnotesize\textcolor{doccolor}{\textbf{Doc~3}}}\hfill{\tiny\texttt{score:~26.089}}\ {\tiny\textcolor{gray}{\{'relevant': \{'868c6e58-57c1-48ee-8c67-832dbedcac18': \{'value': 'no', 'timestamp': 1724071418\}\}\}}}\\[3pt]
{\small\textit{(sans titre)}}\\[2pt]
\vspace{2pt}
{\footnotesize 
"EBITDA is in my opinion not a useful measure for an investor looking to buy shares on the stock market. It is more useful for private businesses open to changing their structuring, or looking to sell significant parts of their business. One of the main benefits of reporting Earnings Before Interest, Taxes, Depreciation \& Amortization, is that it presents the company as it would look to a potential buyer. Consider that net income, as a metric, includes interest costs, taxes, and depreciation. Interest costs are (to put it simply) a result of multiplying a business's debt by its interest rate. If you own a business, and personally guarantee the loan that the company has with the bank, your interest rates might be artificially low. If you have a policy of reaching high debt levels relative to your equity, in order to achieve high 'financial leveraging', your interest cost might be artificially high. Either way, if I bought your business, my debt structure could be completely different, and therefore your interest costs are not particularly relevant to me, a potential buyer. Instead, I should attempt to anticipate what my own interest costs would be, under my plans for your business. Taxes are a result of many factors, including the corporate structure of the business. If you run your business as a sole proprietorship (ie: no corporation), but I want to buy it under my corporation, then my tax rates could look nothing like yours. Or if we operated in multiple jurisdictions. etc. etc. Instead of using your taxes as an estimate for mine, I should anticipate my taxes based on my plans for your business. Depreciation / amortization is a measure that estimates how much of a business's ""fixed assets"" were ""used up"" during the year. ie: how much wear and tear occurred on your fleet of trucks? It is generally calculated as a \% of your overall asset value. It is a (very loose) proxy for the cash costs which will ultimately be incurred to make repairs/replacements. D\&A is also something which could significantly change if a business changes hands.}
\end{docenv}
\begin{docenv}
{\footnotesize\textcolor{doccolor}{\textbf{Doc~4}}}\hfill{\tiny\texttt{score:~17.648}}\ {\tiny\textcolor{gray}{\{'relevant': \{'868c6e58-57c1-48ee-8c67-832dbedcac18': \{'value': 'yes', 'timestamp': 1724071728\}\}\}}}\\[3pt]
{\small\textit{(sans titre)}}\\[2pt]
\vspace{2pt}
{\footnotesize 
Stated plainly... it's a benefit. Companies are not required to offer you any compensation above paying you minimum wage. But benefits attract higher quality employees.  I think a big part of it is that it is the norm. Employees want it because of the tax benefits. Employees expect it because almost all reputable companies of any significant size offer it. You could run a great company, but if you don't offer a 401k plan, you can scare away good potential employees. It would give a bad impression the same way that not offering health insurance would.}
\end{docenv}
\begin{docenv}
{\footnotesize\textcolor{doccolor}{\textbf{Doc~5}}}\hfill{\tiny\texttt{score:~15.389}}\ {\tiny\textcolor{gray}{\{'relevant': \{'868c6e58-57c1-48ee-8c67-832dbedcac18': \{'value': 'yes', 'timestamp': 1724071941\}\}\}}}\\[3pt]
{\small\textit{(sans titre)}}\\[2pt]
\vspace{2pt}
{\footnotesize 
The benefits Apple offers are pretty amazing. 15\% discount on stock (you can allocate up to 10\% of your paycheck to stock purchases). 25\% 401K matching (apple matches 25\%. You can allocate up to 25\% of your paycheck to the company 401k). Excellent health/vision/dental/etc benefits to anyone working 30 hours a week or more. Employees are usually given an iPhone or iPod to have, so that they are seen by customers using Apple products.  Best Buy and Staples probably don't do all of these things for their salespeople.}
\end{docenv}
\begin{docenv}
{\footnotesize\textcolor{doccolor}{\textbf{Doc~6}}}\hfill{\tiny\texttt{score:~10.750}}\ {\tiny\textcolor{gray}{\{'relevant': \{'868c6e58-57c1-48ee-8c67-832dbedcac18': \{'value': 'yes', 'timestamp': 1724072268\}\}\}}}\\[3pt]
{\small\textit{(sans titre)}}\\[2pt]
\vspace{2pt}
{\footnotesize 
The majority of the employers hold the misconception that employee training programs are costly and an unwanted expense. But the fact is that a company can reap plenty of benefits by providing personnel training to their employees. Here below we can find some of the top benefits one could enjoy from an employee training program.}
\end{docenv}
\begin{refanswerenv}
{\footnotesize\textit{\textcolor{gray}{\textbf{Reference answer:}\enspace There are several benefits that a company may offer to its employees in addition to their salary. These can include:
1. Health Insurance: This can include medical, dental, and vision coverage. Some companies may cover the entire cost of premiums, while others may require employees to contribute a portion.
2. Retirement Benefits: Companies may offer retirement plans such as a 401(k) or pension plan. They may also match a certain percentage of employee contributions.
3. Paid Time Off: This can include vacation days, sick days, personal days, and holidays.
4. Stock Options: Some companies may offer employees the opportunity to purchase company stock at a discounted price.
5. Professional Development: Companies may offer tuition reimbursement for further education or training programs related to the employee's job.
6. Childcare Assistance: Some companies may offer childcare subsidies or on-site childcare facilities.
7. Flexible Working Hours: Some companies may offer flexible working hours or remote work options to accommodate employees' personal needs.
8. Performance Bonuses: Companies may offer bonuses based on individual or team performance.
9. Employee Discounts: Companies may offer discounts on their products or services, or discounts at other businesses.
10. Wellness Programs: Some companies may offer wellness programs that include gym memberships, health screenings, or stress management resources.}}}
\end{refanswerenv}
\vspace{0.5em}
\noindent{\small\textcolor{usercolor}{\textbf{$\triangleright$\enspace User --- Turn~7}}}\hfill{\tiny\textcolor{gray}{Factoid, Summarization $\cdot$ Follow-up $\cdot$ ANSWERABLE}}
\begin{userbubble}
By the way, what are the expenses that a company has with its employees besides their remuneration?
\end{userbubble}
\noindent\hfill{\small\textcolor{agentcolor}{\textbf{Agent\enspace$\triangleleft$}}}
\begin{agentbubble}
Apart from employee remuneration, companies may incur various other expenses related to their employees as taxes, benefits and capital. There are several taxes that an employer has to pay  for hiring someone, including social security taxes (which goes to the employee retirement), unemployment insurance tax (the unemployment benefits if they lay off the employee), and workers compensation tax (pays if the workers are injured on the job); Benefits: In the U.S. employers often pay for medical insurance, sometimes for dental, life, or disability. There's usually some sort of retirement plan. They expect to give you some vacation days, holidays, and sick days where they pay you even though you're not working.  Capital: Often the company has to provide you with some sort of equipment, like a computer; furniture, like a chair and desk, etc. Some companies also offer their staff a day care center and that is a considerably large business expense that a lower margin business could no afford. 
\end{agentbubble}
\noindent\hfill{\footnotesize\textcolor{doccolor}{\textbf{$\hookrightarrow$\enspace Retrieved documents~(5)}}}
\begin{docenv}
{\footnotesize\textcolor{doccolor}{\textbf{Doc~1}}}\hfill{\tiny\texttt{score:~42.686}}\ {\tiny\textcolor{gray}{\{'relevant': \{'868c6e58-57c1-48ee-8c67-832dbedcac18': \{'value': 'no', 'timestamp': 1724074028\}\}\}}}\\[3pt]
{\small\textit{(sans titre)}}\\[2pt]
\vspace{2pt}
{\footnotesize 
I don't have any experience in this, but this is my academic understanding of business pricing.  The LOWEST amount a seller would accept is the liquidation value. For a B\&B, what would the value of the land, the house, the furnishings, accounts payable, etc. be if it had to be sold today, minus any liabilities.  The amount the seller would like to pay for is going to be a multiple of its annual earnings. One example of this is the discounted cash flow analysis. You determine the EBITDA, the earnings a company generated, before interest,  depreciation, taxation and amortization.  Once you have this amount, you can project it out in perpetuity, or you use an industry multiplier.  Perpetuity:  You project this value out in perpituity, discounted by the going interest rate. In other words, if you project the business will earn \$100,000/year, the business should grow at a 5\% rate, and the going interest rate is 8\%. Using a growing perpetuity formula, one value of a business would be: 100,000 / (.08 - .03) = \$2,000,000.  This is a very high number, and the seller would love to get it.  It's more common to do a multiple of the EBIDTA. You can do some research into the valuation of the particular industry to figure out the EBIDTA multiplier for the industry. For example, this article suggests that the 2011 EBITDA multiplier for hospitality industries is 13.8. (It's valuing large hotel chains, but it's a start). So the value of this B\&B would be around \$1,380,000. Here is an online SME valuation tool to help with the EBIDTA multiple based valuation. Also, from my research, it looks like many small business use Seller Discretionary Earnings  (SDE) instead of EBITDA. I don't know much about it, but it seems to serve a similar purpose as EBITDA.  A potential buyer should request the financial statements of the business for the last few years to determine the value of the business, and then can negotiate with the owner a price. You would probably want to enlist a broker to help you with the transaction.}
\end{docenv}
\begin{docenv}
{\footnotesize\textcolor{doccolor}{\textbf{Doc~2}}}\hfill{\tiny\texttt{score:~30.902}}\ {\tiny\textcolor{gray}{\{'relevant': \{'868c6e58-57c1-48ee-8c67-832dbedcac18': \{'value': 'no', 'timestamp': 1724073313\}\}\}}}\\[3pt]
{\small\textit{(sans titre)}}\\[2pt]
\vspace{2pt}
{\footnotesize 
You need experience. Work in retail, hopefully to some sort of managerial level where you get to see behind the scenes and not just manage associates.  If you're really gung-ho about it and willing to take a big risk, I've heard many people say the same reason that a small business fails - not enough cash flow. You might be lucky to raise enough capital for start-up costs, but can you survive two years of negative income? And I'm not even talking about your personal living expenses. Forget about that, you're going to be living at home eating ramen. I'm talking about having enough cash to pay the rent and keep the lights on. Unfortunately, it's going to be a vicious cycle because banks know small businesses are a huge risk and are unwilling to provide loans for them to survive...but small businesses fail because they aren't able to survive in the short-term no matter how good the prospects are in the long-term.   Your safest bet is to grow organically. This means having a business model with negligible start-up costs and little overhead costs. Your inventory is going to be small and because of that, your margins are going to be thin (no wholesale prices). But at this point, based on the products you're selling, you're basically just another one of the millions of internet re-sellers. You're going to basically have to place all your value in getting savvy to learn how to spot good deals and re-sell products for a higher price.   Honestly, I'll be harsh about it. You have 0 competitive advantage and offer 0 extra value in your idea. There's really no point in wasting your time and money unless you can pivot the idea so you're creating actual competitive advantages for yourself. Otherwise, just stick to re-selling online, at the very least to get a feel for the profit margins and gain experience.}
\end{docenv}
\begin{docenv}
{\footnotesize\textcolor{doccolor}{\textbf{Doc~3}}}\hfill{\tiny\texttt{score:~26.089}}\ {\tiny\textcolor{gray}{\{'relevant': \{'868c6e58-57c1-48ee-8c67-832dbedcac18': \{'value': 'no', 'timestamp': 1724073315\}\}\}}}\\[3pt]
{\small\textit{(sans titre)}}\\[2pt]
\vspace{2pt}
{\footnotesize 
"EBITDA is in my opinion not a useful measure for an investor looking to buy shares on the stock market. It is more useful for private businesses open to changing their structuring, or looking to sell significant parts of their business. One of the main benefits of reporting Earnings Before Interest, Taxes, Depreciation \& Amortization, is that it presents the company as it would look to a potential buyer. Consider that net income, as a metric, includes interest costs, taxes, and depreciation. Interest costs are (to put it simply) a result of multiplying a business's debt by its interest rate. If you own a business, and personally guarantee the loan that the company has with the bank, your interest rates might be artificially low. If you have a policy of reaching high debt levels relative to your equity, in order to achieve high 'financial leveraging', your interest cost might be artificially high. Either way, if I bought your business, my debt structure could be completely different, and therefore your interest costs are not particularly relevant to me, a potential buyer. Instead, I should attempt to anticipate what my own interest costs would be, under my plans for your business. Taxes are a result of many factors, including the corporate structure of the business. If you run your business as a sole proprietorship (ie: no corporation), but I want to buy it under my corporation, then my tax rates could look nothing like yours. Or if we operated in multiple jurisdictions. etc. etc. Instead of using your taxes as an estimate for mine, I should anticipate my taxes based on my plans for your business. Depreciation / amortization is a measure that estimates how much of a business's ""fixed assets"" were ""used up"" during the year. ie: how much wear and tear occurred on your fleet of trucks? It is generally calculated as a \% of your overall asset value. It is a (very loose) proxy for the cash costs which will ultimately be incurred to make repairs/replacements. D\&A is also something which could significantly change if a business changes hands.}
\end{docenv}
\begin{docenv}
{\footnotesize\textcolor{doccolor}{\textbf{Doc~4}}}\hfill{\tiny\texttt{score:~17.340}}\ {\tiny\textcolor{gray}{\{'relevant': \{'868c6e58-57c1-48ee-8c67-832dbedcac18': \{'value': 'yes', 'timestamp': 1724073997\}\}\}}}\\[3pt]
{\small\textit{(sans titre)}}\\[2pt]
\vspace{2pt}
{\footnotesize 
(1). Is this right? Pretty much, though this is a really rudimentary way to think about it. (2). If it is, why is it that extensive services are provided by high margin companies competing for talent, rather then lower margin businesses looking to boost their profits by reducing their expenditures on employees (by cutting out the government)? It's the polar opposite of that. Google (and companies like that) do things like have a day care center on premises. The company staffs a day care center which has costs, then lets employees use it for free. This is a business expense for Google, and in relative terms, a considerably large business expense that a lower margin business could no afford. Employer healthcare is a tax protected expense for employees via section 125 of the tax code. The company portion of the healthcare costs are a deductible business expense to the company, as expected. Healthcare is different than most other expenses because the employee can forego income before it's effectively received which negates it from taxable income. This doesn't work for something like food purchased at a cafe on a Google complex. If employee money is being spent at a corporate cafe, it's taxable income being spent (though the cost of running the cafe is a tax deductible business expense to the company). There have been discussions in congress to assess a value as income to employees for services like on site child care and no cost employee cafeterias. To address your new example: For example, suppose John Doe makes \$100,000 a year taxed at a rate of 20\%, for a take home pay of \$80,000. He spends \$10,000 on food. His employer Corporation decides to give him all of his food and deduct it as a business expense - costing them \$10,000. But now they can pay John Doe an amount so his take home pay will be reduced by \$10,000 - \$87,500 The company is now spending \$97500 employing John Doe, for a savings of \$2500\$. This would be an audit prone administrative nightmare.}
\end{docenv}
\begin{docenv}
{\footnotesize\textcolor{doccolor}{\textbf{Doc~5}}}\hfill{\tiny\texttt{score:~14.700}}\ {\tiny\textcolor{gray}{\{'relevant': \{'868c6e58-57c1-48ee-8c67-832dbedcac18': \{'value': 'yes', 'timestamp': 1724073703\}\}\}}}\\[3pt]
{\small\textit{(sans titre)}}\\[2pt]
\vspace{2pt}
{\footnotesize 
"An employee costs the company in four ways: Salary, taxes, benefits, and capital. Salary: The obvious one, what they pay you. Taxes: There are several taxes that an employer has to pay for the privilege of hiring someone, including social security taxes (which goes to your retirement), unemployment insurance tax (your unemployment benefits if they lay you off), and workers compensation tax (pays if you are injured on the job). (There may be other taxes that I'm not thinking of, but in any case those are the main ones.) Benefits: In the U.S. employers often pay for medical insurance, sometimes for dental, life, and disability. There's usually some sort of retirement plan. They expect to give you some number of vacation days, holidays, and sick days where they pay you even though you're not working. Companies sometimes offer other benefits, like discounts on buying company products, membership in health clubs, etc. Capital: Often the company has to provide you with some sort of equipment, like a computer; furniture, like a chair and desk; etc. As far as the company is concerned, all of the above are part of the cost of having you as an employee. If they would pay a domestic employee \$60,000 in salary and \$20,000 in taxes, then assuming the same benefits and capital investment, if a foreign employee would cost them \$0 in taxes they should logically be willing to pay \$80,000. Any big company will have accountants who figure out the total cost of a new employee in excruciating detail, and they will likely be totally rational about this. A smaller company might think, ""well, taxes don't really count ..."" This is irrational but people are not always rational. I don't know what benefits they are offering you, if any, and what equipment they will provide you with, if any. I also don't know what taxes, if any, a U.S. company has to pay when hiring a remote employee in a foreign country. If anybody on here knows the answer to that, please chime in. Balanced against that, the company likely sees disadvantages to hiring a foreign remote employee, too. Communication will be more difficult, which may result in inefficiency.}
\end{docenv}
\begin{refanswerenv}
{\footnotesize\textit{\textcolor{gray}{\textbf{Reference answer:}\enspace Apart from employee remuneration, companies may incur various other expenses related to their employees. These can include costs for benefits such as health insurance, dental insurance, vision care, retirement savings plans (like 401k matching), and stock options. Additionally, there can be costs associated with training and developing employees, providing them with necessary equipment or tools, and offering perks like flexible working hours or remote work options. Other expenses might include taxes, workers' compensation insurance, and costs related to recruiting and hiring new employees. It's also important to consider overhead costs, such as office space, utilities, and maintenance, which can indirectly be linked to employees.}}}
\end{refanswerenv}
\end{document}
